% =====================================================================
%  30.1  The Pruefer bijection.
%  Exports: sec:c32-prufer  def:c32-ltree  def:c32-code  def:c32-dec
%  thm:c32-bij  lem:c32-deg  cor:c32-mult  lst:c32-encode
%  lst:c32-decode  fig:c32-encode  fig:c32-decode
%  This section OWNS the Pruefer encode/decode pair, the degree-from-code
%  lemma, and the degree-multiset multinomial count; 30.2-30.5 and the
%  Summit reference them by label.  NO solved-problem block here (the
%  solved problems are 2190D in 30.2 and 280C in 30.4).  Everything is
%  worked examples on the running 6-vertex tree T with code [5,5,3,4].
% =====================================================================
\section{The Pr\"ufer Bijection}\label{sec:c32-prufer}

\begin{previouslybox}[Ch.29 extracted one tree from a graph]
Chapter~29 (\chref{c31}) ended with a graph and the tree that hides inside
it -- a bridge tree, a block-cut tree, a rho decomposition.  Chapter~30 opens
with the opposite question: the input is a \emph{description}, and the answer
is a number.  To count trees you need a way to walk through \(\mathcal{T}_n\)
without enumerating it, and that is what the Pr\"ufer bijection is -- a
canonical code, one sequence per labelled tree, in a form you can count with
elementary combinatorics.
\end{previouslybox}

\begin{topicintro}
A \term{labelled tree} on \([n]\) is a tree whose vertices are the specific
set \(\{1,2,\dots,n\}\) (\(\defref{c32-ltree}\)); two labelled trees differ
when their edge sets do, so a relabelling that moves an edge is a \emph{new}
tree.  Pr\"ufer's construction (\(\defref{c32-code}\)) is an iterative peel:
as long as more than two vertices remain, find the leaf with the smallest
label, delete it, and append the label of its (unique) neighbour to a running
list.  After \(n-2\) steps the list is the \term{Pr\"ufer code} of \(T\), a
sequence of length \(n-2\) with entries in \([n]\).  The inverse
(\(\defref{c32-dec}\)) reconstructs \(T\) from its code greedily: at each step
take the smallest label of \([n]\) that is not in the residual suffix of the
code, join it to the code's first remaining entry, and drop that entry.  The
two maps are mutually inverse (\(\thmref{c32-bij}\)), which is the whole point
of the chapter: counting labelled trees is counting sequences.

The most-used corollary is \(\lemref{c32-deg}\): the number of times a label
\(v\) appears in \(\mathrm{code}(T)\) is \(\deg_T(v)-1\).  Read a degree
constraint as a multiplicity constraint on the code, and the count becomes a
multinomial (\(\corref{c32-mult}\)).

\textbf{Prerequisites.} The notion of a leaf (a vertex of degree 1) and the
handshake identity \(\sum_v \deg(v) = 2(n-1)\), both from Ch.~1.  Nothing
about trees beyond "connected and \(n-1\) edges" is assumed here; the
bijection itself is proved from scratch.

\textbf{After this section you can}
\begin{itemize}[nosep]
  \item write the Pr\"ufer code of any labelled tree in \(\bigO{n\log n}\) via a min-heap of leaves; \hfill\lv{3}
  \item decode a length-\((n-2)\) sequence over \([n]\) back to its tree in \(\bigO{n\log n}\); \hfill\lv{3}
  \item state and use \(\lemref{c32-deg}\) to convert degree constraints into multiplicity constraints; \hfill\lv{4}
  \item count labelled trees with a prescribed degree sequence via \(\corref{c32-mult}\). \hfill\lv{4}
\end{itemize}
\end{topicintro}

% ---------------------------------------------------------------
\subsection{Encoding a labelled tree to its Pr\"ufer code}
% ---------------------------------------------------------------

\begin{definition}{Labelled tree on \([n]\)}{c32-ltree}
A \term{labelled tree} on \([n]\) is an undirected graph \(T=(V,E)\) with
\(V=[n]\), connected and with \(|E|=n-1\).  The set of all such trees is
\(\mathcal{T}_n\).  Two labelled trees \(T,T'\) are equal iff they have the
same edge set \(\{u,v\}\subseteq[n]\); a permutation of the labels produces a
\emph{different} tree unless it fixes every edge.
\end{definition}

\begin{definition}{Pr\"ufer code (encoding)}{c32-code}
Given \(T\in\mathcal{T}_n\) with \(n\ge 2\), form a list \(c\) initially
empty.  While \(|V(T)|>2\): let \(\ell\) be the smallest leaf of the current
\(T\); append the label of \(\ell\)'s unique neighbour to \(c\); delete \(\ell\)
and its incident edge.  The final \(c\) is the \term{Pr\"ufer code} of \(T\),
written \(\mathrm{code}(T)\).  It has length exactly \(n-2\), because one
vertex is removed at each of the \(n-2\) iterations.
\end{definition}

\begin{lemma}{Degree from multiplicity}{c32-deg}
For \(T\in\mathcal{T}_n\) and any \(v\in[n]\),
\[
  m_v(T) \;:=\; \#\{\text{positions } i : c_i = v\} \;=\; \deg_T(v) - 1 .
\]
In particular, a vertex is a leaf iff its label never appears in the code.
\end{lemma}

\begin{proofbox}
Track the current degree of each vertex as the peel runs.  When a leaf
\(\ell\) is deleted, we write its neighbour \(v\), and \(v\)'s current degree
drops by one (its edge to \(\ell\) is gone).  Thus \(v\) is written down
exactly at those iterations in which \(v\) survives but one of its incident
edges is being removed by a leaf peel; \(v\) itself is removed only when its
current degree reaches \(1\), and the last-removed vertex is the label that
would have been written at a hypothetical step \(n-1\).  Over the whole
run, the number of edges incident to \(v\) is \(\deg_T(v)\), of which all but
the last one (the edge on which \(v\) itself becomes a leaf, or is retained
at the very end) contribute exactly one write.  Hence \(m_v(T)=\deg_T(v)-1\).
\(\Box\)
\end{proofbox}

\begin{theorem}{Pr\"ufer's bijection}{c32-bij}
For every \(n\ge 2\), the map \(\mathrm{code}:\mathcal{T}_n\to[n]^{n-2}\)
sending a labelled tree to its Pr\"ufer code is a bijection, with inverse
\(\mathrm{decode}\) (\(\defref{c32-dec}\)).
\end{theorem}

\begin{corollary}{Degree-prescribed count}{c32-mult}
Let \(d_1,\dots,d_n\ge 1\) with \(\sum_i d_i=2(n-1)\).  The number of
labelled trees on \([n]\) with \(\deg(i)=d_i\) is
\[
  (n-2)!\,\Big/\,\prod_{i=1}^{n} (d_i-1)! \;=\; \binom{n-2}{d_1-1,\;\dots,\;d_n-1} .
\]
\end{corollary}

\begin{proofbox}
By \(\lemref{c32-deg}\), the degree constraint is exactly the multiplicity
constraint "\(i\) appears \(d_i-1\) times in the code".  A length-\((n-2)\)
sequence over \([n]\) with those multiplicities is a multinomial arrangement;
\(\thmref{c32-bij}\) turns each sequence into a unique tree. \(\Box\)
\end{proofbox}

\begin{definition}{Pr\"ufer decoding}{c32-dec}
Given a length-\((n-2)\) sequence \(c=(c_1,\dots,c_{n-2})\) over \([n]\),
reconstruct \(T\in\mathcal{T}_n\) as follows.  Maintain the residual code
suffix \(c^{(t)}=(c_t,\dots,c_{n-2})\) and the vertex set \([n]\).  At step
\(t=1,\dots,n-2\): let \(\ell_t\) be the smallest label of \([n]\) that does
not appear in \(c^{(t)}\); emit edge \(\{\ell_t,c_t\}\); conceptually delete
\(\ell_t\) from the candidate pool (it cannot be a leaf again) and drop \(c_t\)
from the residual suffix.  After \(n-2\) steps two labels \(a,b\) remain; emit
\(\{a,b\}\).  The result is \(\mathrm{decode}(c)\).
\end{definition}

\begin{lstlisting}[caption={Pruefer encode. Maintain a min-heap of current leaves and a degree array; pop the smallest leaf, write its live neighbour, decrement that neighbour, and push it if it becomes a leaf. ASCII only, identifiers without underscores}, label={lst:c32-encode}]
#include <bits/stdc++.h>
using namespace std;
int main(){
  int n; if(scanf("%d",&n)!=1) return 0;
  vector<vector<int> > adj(n+1);
  vector<int> deg(n+1,0);
  for(int i=0;i<n-1;i++){
    int u,v; scanf("%d %d",&u,&v);
    adj[u].push_back(v); adj[v].push_back(u);
    deg[u]++; deg[v]++;
  }
  priority_queue<int, vector<int>, greater<int> > leaves;
  vector<int> alive(n+1,1);
  for(int v=1;v<=n;v++) if(deg[v]==1) leaves.push(v);
  vector<int> code;
  for(int step=0; step<n-2; step++){
    int l=leaves.top(); leaves.pop();
    int nb=-1;
    for(int x : adj[l]) if(alive[x]){ nb=x; break; }
    code.push_back(nb);
    alive[l]=0; deg[nb]--;
    if(deg[nb]==1) leaves.push(nb);
  }
  for(size_t i=0;i<code.size();i++) printf("%d ", code[i]);
  printf("\n");
  return 0;
}
\end{lstlisting}

\subsection*{Worked Example: Encoding the Running Tree}

The running tree \(T\) from \(\figref{c32-tree}\) has edges
\(\{1\!-\!5,\,2\!-\!5,\,5\!-\!3,\,3\!-\!4,\,4\!-\!6\}\).  Degrees are
\((\deg(1),\dots,\deg(6))=(1,1,2,2,3,1)\), leaves are \(\{1,2,6\}\).  The
peel runs:

\begin{itemize}[nosep]
  \item Step 1: smallest leaf is \(1\); its neighbour is \(5\).  Append \(5\).
        Delete edge \(1\!-\!5\).
  \item Step 2: current leaves are \(\{2,6\}\); smallest is \(2\); neighbour
        is \(5\).  Append \(5\).  Delete edge \(2\!-\!5\).
  \item Step 3: current leaves are \(\{5,6\}\); smallest is \(5\); neighbour
        is \(3\).  Append \(3\).  Delete edge \(5\!-\!3\).
  \item Step 4: current leaves are \(\{3,6\}\); smallest is \(3\); neighbour
        is \(4\).  Append \(4\).  Delete edge \(3\!-\!4\).  Two vertices
        \(\{4,6\}\) remain; stop.
\end{itemize}

The code is \([5,5,3,4]\), exactly as advertised.  Multiplicity check via
\(\lemref{c32-deg}\): \(5\) appears twice \(\Rightarrow \deg(5)=3\); \(3,4\)
appear once \(\Rightarrow\) degree \(2\); \(1,2,6\) appear zero \(\Rightarrow\)
degree \(1\).  The three leaves in the code's zero-set are precisely the
labels that never appear.

\begin{figure}[H]
\centering
\begin{tikzpicture}[diagram, scale=0.85, transform shape,
   tnode/.style={circle, draw=primaryblue!70, fill=primaryblue!8, line width=0.7pt,
                 minimum size=7mm, font=\scriptsize\bfseries, inner sep=1pt},
   gone/.style={circle, draw=gray!50, fill=gray!10, line width=0.5pt,
                 minimum size=7mm, font=\scriptsize, inner sep=1pt}]
  \node[tnode] (a5) at (0,0.9) {5};
  \node[tnode] (a3) at (2,0.9) {3};
  \node[tnode] (a4) at (3.6,0.9) {4};
  \node[tnode] (a6) at (5,0.9) {6};
  \node[gone]  (x1) at (-1.4,-0.6) {1};
  \node[gone]  (x2) at (-0.4,-1.4) {2};
  \draw[primaryblue!80, line width=1.0pt] (a5)--(a3)--(a4)--(a6);
  \draw[gray!60, dashed, line width=0.7pt] (x1)--(a5);
  \draw[gray!60, dashed, line width=0.7pt] (x2)--(a5);
  \node[annot, font=\scriptsize, anchor=north, text width=9.5cm] at (2,-1.6)
    {after dropping \(1,2\): the peel writes \(5,5\); then \(3\) and \(4\) on the shrinking spine};
\end{tikzpicture}
\caption{Mid-encode state after steps 1--2: leaves \(1,2\) dropped, code so
far \([5,5]\), remaining tree is the path \(5\!-\!3\!-\!4\!-\!6\).  Steps 3--4
drop \(5\) (write \(3\)) then \(3\) (write \(4\)).}
\label{fig:c32-encode}
\end{figure}

% ---------------------------------------------------------------
\subsection{Decoding a Pr\"ufer code back to its tree}
% ---------------------------------------------------------------

\begin{lstlisting}[caption={Pruefer decode. Precompute count of each label in the code; use a min-heap of zero-count vertices as the candidate leaves; at each step emit edge smallest-leaf to code at t, decrement count of code at t, and push code at t into the heap if it just hit zero. After n-2 steps, two labels remain, emit their edge. ASCII only, identifiers without underscores}, label={lst:c32-decode}]
#include <bits/stdc++.h>
using namespace std;
int main(){
  int n; if(scanf("%d",&n)!=1) return 0;
  vector<int> code(n-2);
  vector<int> cnt(n+1,0);
  for(int i=0;i<n-2;i++){ scanf("%d",&code[i]); cnt[code[i]]++; }
  priority_queue<int, vector<int>, greater<int> > heap;
  for(int v=1;v<=n;v++) if(cnt[v]==0) heap.push(v);
  vector<pair<int,int> > edges;
  for(int t=0;t<n-2;t++){
    int l=heap.top(); heap.pop();
    edges.push_back({l, code[t]});
    cnt[code[t]]--;
    if(cnt[code[t]]==0) heap.push(code[t]);
  }
  int a=heap.top(); heap.pop();
  int b=heap.top();
  edges.push_back({a,b});
  for(auto &e : edges) printf("%d %d\n", e.first, e.second);
  return 0;
}
\end{lstlisting}

\subsection*{Worked Example: Decoding \([5,5,3,4]\) on \(n=6\)}

Precompute multiplicities: \(\mathrm{cnt} = (0,0,1,1,2,0)\) for labels
\(1,\dots,6\).  Initial heap (zero-count labels) = \(\{1,2,6\}\).

\begin{itemize}[nosep]
  \item \(t=1\): pop \(1\), emit \(1\!-\!5\); \(\mathrm{cnt}[5]\to 1\), not
        zero, do not push.  Heap \(=\{2,6\}\).
  \item \(t=2\): pop \(2\), emit \(2\!-\!5\); \(\mathrm{cnt}[5]\to 0\),
        push \(5\).  Heap \(=\{5,6\}\).
  \item \(t=3\): pop \(5\), emit \(5\!-\!3\); \(\mathrm{cnt}[3]\to 0\),
        push \(3\).  Heap \(=\{3,6\}\).
  \item \(t=4\): pop \(3\), emit \(3\!-\!4\); \(\mathrm{cnt}[4]\to 0\),
        push \(4\).  Heap \(=\{4,6\}\).
  \item Final pair \(\{4,6\}\), emit \(4\!-\!6\).
\end{itemize}

The recovered edges are \(\{1\!-\!5,\,2\!-\!5,\,5\!-\!3,\,3\!-\!4,\,4\!-\!6\}\)
-- exactly the running tree \(T\).

\begin{figure}[H]
\centering
\begin{tikzpicture}[diagram, scale=0.9, transform shape,
   box/.style={rectangle, draw=primaryblue!70, fill=primaryblue!8,
               minimum width=8mm, minimum height=7mm, font=\scriptsize\bfseries, inner sep=1pt},
   arr/.style={-Latex, line width=0.9pt, draw=primaryblue!70}]
  \node[box] (c1) at (0,1.6) {5};
  \node[box] (c2) at (1.2,1.6) {5};
  \node[box] (c3) at (2.4,1.6) {3};
  \node[box] (c4) at (3.6,1.6) {4};
  \node[box] (l1) at (0,0) {1};
  \node[box] (l2) at (1.2,0) {2};
  \node[box] (l3) at (2.4,0) {5};
  \node[box] (l4) at (3.6,0) {3};
  \node[box] (l5) at (4.8,0) {4 -- 6};
  \draw[arr] (c1)--(l1);
  \draw[arr] (c2)--(l2);
  \draw[arr] (c3)--(l3);
  \draw[arr] (c4)--(l4);
  \node[annot, font=\scriptsize, anchor=north] at (2.4,-1.4)
    {at each \(t\), the smallest label absent from the residual suffix is the leaf};
\end{tikzpicture}
\caption{Decode of \([5,5,3,4]\): the four code entries generate four edges
one-for-one (\(1\!-\!5, 2\!-\!5, 5\!-\!3, 3\!-\!4\)); the final residual pair
\(\{4,6\}\) supplies the fifth edge \(4\!-\!6\).}
\label{fig:c32-decode}
\end{figure}

\begin{insightbox}[Count trees with a fixed degree sequence]
Suppose \(n=5\) and the desired degrees are \((\deg(1),\dots,\deg(5)) =
(3,1,1,1,2)\).  Feasibility check: \(\sum d_i = 8 = 2(n-1)\) passes.
By \(\lemref{c32-deg}\) the code multiplicities are \(m_1=2,\, m_2=m_3=m_4=0,
\,m_5=1\), so the multinomial \(\corref{c32-mult}\) gives \((n-2)!/\prod m_i!
= 3!/(2!\cdot 1!) = 3\) labelled trees.  Their three codes are the three
permutations of the multiset \(\{5,1,1\}\): \([5,1,1], [1,5,1], [1,1,5]\).
Decoding \([5,1,1]\) via \(\coderef{c32-decode}\): initial \(\mathrm{cnt} =
(2,0,0,0,1)\); initial zero-count heap \(= \{2,3,4\}\).  Step 1 pops \(2\),
suffix is \([1,1]\), labels absent from suffix are \(\{2,3,4,5\}\), emit
\(2\!-\!5\) and push \(5\) into the heap once \(\mathrm{cnt}[5]\to 0\).
Step 2 pops \(3\) (smallest still absent from residual \([1]\) among unused),
emit \(3\!-\!1\), \(\mathrm{cnt}[1]\to 1\).  Step 3 pops \(4\), emit \(4\!-\!1\),
\(\mathrm{cnt}[1]\to 0\), push \(1\).  Two labels \(\{1,5\}\) remain; emit
\(1\!-\!5\).  Result: edges \(\{2\!-\!5,\, 3\!-\!1,\, 4\!-\!1,\, 1\!-\!5\}\),
degrees \((3,1,1,1,2)\) as required.  The other two codes yield two further
non-isomorphic labelings; total \(= 3\), matching the multinomial.
\end{insightbox}

\begin{insightbox}[Why the bijection is what the chapter is built on]
Every "count trees with property \(P\)" question in this chapter reduces to
"count sequences with property \(P'\)", where \(P'\) is what \(P\) becomes in
code space.  Degree constraints become multiplicity constraints
(\(\corref{c32-mult}\)).  Structural constraints -- "connected to \(v\)",
"leaf" -- become statements about the code's zero-set or the position of the
label at peel time.  This is not merely a proof trick: it is \emph{the}
algorithm you implement (CF 2190D \(\secref{c32-cayley}\) and the Summit
\(\probref{c32-summit}\) both run an encode/decode-style loop under a
counting filter, not a DFS).
\end{insightbox}
